% Example 1: a Figure 1 teaser with three parallel points, one per card, for a fictitious benchmark on
% repairing a repository from an issue report. The cards follow the task left to right: the input (an issue),
% the output (a patch) and the judgement (hidden tests). Numbered badges match an enumeration in the text.
% Orange marks the open decision in each card; the pill states it in one line; the footer band gives the baseline.
% Build:  python <skill>/scripts/build_figure.py teaser-points.tex
\documentclass[10pt,tikz,border={1pt 3.5pt 3pt 1.5pt}]{standalone}
\usepackage{cardfig}
\setlength{\cardw}{4.35cm}
\setlength{\cardh}{4.45cm}
\setlength{\cardgap}{0.35cm}
\begin{document}
\begin{tikzpicture}[remember picture]
\cardpanels{3}
\cardheadn{P1}{1}{Issue}  \cardicon{P1}{bug}
\cardheadn{P2}{2}{Patch}  \cardicon{P2}{code}
\cardheadn{P3}{3}{Tests}  \cardicon{P3}{vial}
\cardfoot{P1}{earlier work: the file to edit is given}
\cardfoot{P2}{earlier work: one candidate per task}
\cardfoot{P3}{earlier work: tests shown to the system}
\cardpill{P1}{the issue names no file}
\cardpill{P2}{many candidates, one is right}
\cardpill{P3}{hidden tests decide}
% ================= card 1: the issue text on the left, the repository files on the right, the file choice open
\docpage[text width=1.55cm]{d}{$(P1.north west)+(0.28cm,-0.95cm)$}{ISSUE}{%
  Empty input makes \tikzmarknode[hl]{m1}{parse()} raise KeyError. Seen since \tikzmarknode[hl]{m2}{v2.3}.}
\node[grp, anchor=east] at ($(P1.north east)+(-0.3cm,-0.92cm)$) {\faIcon{code-branch}\ repository};
\node[colA, anchor=east] at ($(P1.north east)+(-0.3cm,-1.5cm)$) (f1) {parse.py};
\node[colA, anchor=east] at ($(P1.north east)+(-0.3cm,-2.05cm)$) (f2) {io.py};
\node[colA, anchor=east] at ($(P1.north east)+(-0.3cm,-2.6cm)$) (f3) {cli.py};
% the badge sits on the connector between the highlighted phrase and the files; positions from marks need overlay
\coordinate (qx) at ($(P1.north west)+(2.55cm,0)$);
\node[qb, overlay] at (qx |- m1.center) (q1) {?};
\draw[wire, overlay] (d.east |- m1.center) -- (q1);
\draw[arr, overlay] (q1) -- (f1.west);
\draw[arr, overlay] (q1) -- (f2.west);
\draw[arr, overlay] (q1) -- (f3.west);
% ================= card 2: several candidate patches, one accepted
\node[grp, anchor=west] at ($(P2.north west)+(0.3cm,-0.92cm)$) {\faIcon{layer-group}\ candidates};
\node[grp, anchor=east] at ($(P2.north east)+(-0.3cm,-0.92cm)$) {{\color{kept}\faIcon{check}}\ accepted};
\node[colB, anchor=west] at ($(P2.north west)+(0.3cm,-1.4cm)$)  (c1) {patch A};
\node[colB, anchor=west] at ($(P2.north west)+(0.3cm,-1.95cm)$) (c2) {patch B};
\node[colB, anchor=west] at ($(P2.north west)+(0.3cm,-2.5cm)$)  (c3) {patch C};
\node[prop, anchor=east] at ($(P2.north east)+(-0.3cm,-1.95cm)$) (a1) {patch};
\node[qb] at ($(P2.north)+(0,-1.95cm)$) (q2) {?};
\draw[wire] (c1.east) -- (q2);
\draw[wire] (c2.east) -- (q2);
\draw[wire] (c3.east) -- (q2);
\draw[arr] (q2) -- (a1.west);
% ================= card 3: the hidden tests, the one that flips is highlighted
\node[tname, text=dbA, anchor=south west] at ($(P3.north west)+(0.45cm,-1.14cm)$) {\faIcon{eye-slash}\ hidden tests};
\matrix[tblA, anchor=north west,
  column 1/.style={nodes={text width=1.35cm}}, column 2/.style={nodes={text width=0.75cm}}, column 3/.style={nodes={text width=0.75cm}}]
  at ($(P3.north west)+(0.45cm,-1.24cm)$) (tt) {
  test & before & after \\
  test\_empty & fail & |[hi]| pass \\
  test\_parse & pass & pass \\
  test\_cli & pass & pass \\
};
\boxshadow[]{tt}
% the verdict chip under the column that flipped; the arrow leaves the table bottom under that column
\node[prop, anchor=north east] at ($(tt.south east)+(0,-0.35cm)$) (r1) {resolved};
\draw[arr] (tt.south -| tt-2-3) -- (r1.north -| tt-2-3);
\end{tikzpicture}
\end{document}
